\begin{helper}
Let \(G\) be a finite graph, let \(C\subseteq V(G)\) be a finite candidate
set, and let \(Q\subseteq V(G)\) be a finite coordinate set.  Suppose \(Q\)
contains every neighbor of every candidate: whenever \(x\in C\) and
\(xy\in E(G)\), one has \(y\in Q\).  Then, for every activation--priority
outcome \((A,\pi)\), the event \(I_G(A,\pi)\cap C\ne\emptyset\) is equivalent
to the finite local survival event
\[
\exists x\in C\cap A\quad
  \text{such that}\quad
  \pi(y)<\pi(x)
  \text{ for every }y\in A\cap Q\text{ with }xy\in E(G).
\]
\end{helper}
\begin{proof}
By definition of the activation--priority output,
\[
I_G(A,\pi)=\{v\in A:\pi(y)<\pi(v)\hbox{ for every }y\in A\cap N_G(v)\}.
\]
Thus \(I_G(A,\pi)\cap C\ne\emptyset\) holds if and only if there is an
\(x\in C\cap A\) such that \(\pi(y)<\pi(x)\) for every activated neighbor
\(y\) of \(x\).

If such a surviving \(x\) exists, then it certainly beats every activated
neighbor lying in \(Q\), so the finite local survival condition holds.

Conversely, suppose the finite local survival condition holds for some
\(x\in C\cap A\).  Let \(y\in A\) be any neighbor of \(x\).  Since \(x\in C\),
the covering hypothesis places \(y\) in \(Q\).  Therefore \(y\) is one of the
neighbors tested by the finite local condition, and hence
\(\pi(y)<\pi(x)\).  This proves that \(x\) beats every activated neighbor in
the full graph, so \(x\in I_G(A,\pi)\) and
\(I_G(A,\pi)\cap C\ne\emptyset\).  The two events are pointwise identical.
\end{proof}
